\textbf{Phần 1. Trắc nghiệm lựa chọn}
\setcounterpageref{ex}{0}

\begin{ex} %[2D8N38]
Trong phép tính $2 \times 6 = 12$, số 12 được gọi là
\choice
{Thừa số}
{Số bị chia}
{\True Tích}
{Thương}
\loigiai{
Kết quả của phép nhân gọi là Tích.\\
\textbf{Đáp án của câu hỏi là C}
}
\end{ex}

\begin{ex} %[2D8N38]
Trong phép tính $5 \times 4 = 20$, các số 5 và 4 được gọi là
\choice
{Tích}
{\True Thừa số}
{Số hạng}
{Số bị trừ}
\loigiai{
5 và 4 là các thừa số.\\
\textbf{Đáp án của câu hỏi là B}
}
\end{ex}

\begin{ex} %[2D8H38]
Tìm tích biết hai thừa số lần lượt là 2 và 8:
\choice
{\True $16$}
{$14$}
{$18$}
{$10$}
\loigiai{
$2 \times 8 = 16$.\\
\textbf{Đáp án của câu hỏi là A}
}
\end{ex}

\begin{ex} %[2D8H38]
Tìm tích biết hai thừa số lần lượt là 5 và 7:
\choice
{$30$}
{$40$}
{$45$}
{\True $35$}
\loigiai{
$5 \times 7 = 35$.\\
\textbf{Đáp án của câu hỏi là D}
}
\end{ex}

\begin{ex} %[2D8H38]
Khi đổi chỗ hai thừa số trong một tích thì tích sẽ
\choice
{Tăng gấp đôi}
{Giảm đi một nửa}
{\True Không thay đổi}
{Bằng 0}
\loigiai{
Khi đổi chỗ các thừa số thì tích không thay đổi.\\
\textbf{Đáp án của câu hỏi là C}
}
\end{ex}

\begin{ex} %[2D8H38]
Tích của $2 \times 9$ có giá trị bằng tích nào sau đây?
\choice
{$9 + 2$}
{$2 + 9$}
{$2 \times 8$}
{\True $9 \times 2$}
\loigiai{
$2 \times 9 = 9 \times 2 = 18$.\\
\textbf{Đáp án của câu hỏi là D}
}
\end{ex}

\begin{ex} %[2D8V38]
Tích của hai thừa số là 18. Biết một thừa số là 2, thừa số kia là
\choice
{\True $9$}
{$8$}
{$7$}
{$6$}
\loigiai{
$18 : 2 = 9$.\\
\textbf{Đáp án của câu hỏi là A}
}
\end{ex}

\begin{ex} %[2D8V38]
Tích của hai thừa số là 40. Biết một thừa số là 5, thừa số kia là
\choice
{$7$}
{$9$}
{$6$}
{\True $8$}
\loigiai{
$40 : 5 = 8$.\\
\textbf{Đáp án của câu hỏi là D}
}
\end{ex}

\begin{ex} %[2D8V38]
Một thừa số là 5, thừa số kia là số liền sau số 8. Tích của hai số đó là
\choice
{$40$}
{\True $45$}
{$50$}
{$35$}
\loigiai{
Số liền sau của 8 là 9. Tích: $5 \times 9 = 45$.\\
\textbf{Đáp án của câu hỏi là B}
}
\end{ex}

\begin{ex} %[2D8V38]
Một thừa số là 2, thừa số kia là số lớn nhất có một chữ số. Tích là
\choice
{\True $18$}
{$16$}
{$20$}
{$14$}
\loigiai{
Số lớn nhất có 1 chữ số là 9. Tích: $2 \times 9 = 18$.\\
\textbf{Đáp án của câu hỏi là A}
}
\end{ex}

\begin{ex} %[2D8V38]
Tính giá trị biểu thức: $2 \times 8 + 5 \times 4$.
\choice
{$34$}
{$38$}
{\True $36$}
{$40$}
\loigiai{
$16 + 20 = 36$.\\
\textbf{Đáp án của câu hỏi là C}
}
\end{ex}

\begin{ex} %[2D8C38]
Tích của hai số là 10. Nếu giữ nguyên một thừa số và gấp thừa số kia lên 2 lần thì tích mới là
\choice
{$15$}
{\True $20$}
{$25$}
{$30$}
\loigiai{
$10 \times 2 = 20$.\\
\textbf{Đáp án của câu hỏi là B}
}
\end{ex}

\textbf{Phần 2. Trắc nghiệm đúng sai}
\setcounterpageref{ex}{0}

\begin{ex} %[2D8H38]
Xét tính đúng sai về thành phần của phép nhân:
\choiceTF
{\True Trong phép nhân $5 \times 6 = 30$, 5 và 6 là các thừa số.}
{\True 30 được gọi là tích.}
{$5 \times 6 = 6 + 5$.}
{\True $5 \times 6 = 6 \times 5 = 30$.}
\loigiai{
\begin{itemize}
    \item 5 và 6 là thừa số: Đúng.
    \item 30 là tích: Đúng.
    \item $5 \times 6 = 30 \neq 11$. Ý 3 Sai, Ý 4 Đúng.
\end{itemize}
\begin{itemize}
    \item a) \textbf{Đúng}.
    \item b) \textbf{Đúng}.
    \item c) \textbf{Sai}.
    \item d) \textbf{Đúng}.
\end{itemize}
}
\end{ex}

\begin{ex} %[2D8V38]
Cho phép nhân $2 \times 7 = 14$. Xét tính đúng sai:
\choiceTF
{\True Tích là 14.}
{\True Thừa số thứ nhất là 2, thừa số thứ hai là 7.}
{Nếu tăng thừa số 2 thêm 1 đơn vị thì tích tăng thêm 1 đơn vị.}
{\True Khi đó tích là $3 \times 7 = 21$ (tăng thêm 7 đơn vị).}
\loigiai{
\begin{itemize}
    \item Tích là 14: Đúng.
    \item Thừa số là 2 và 7: Đúng.
    \item Tích tăng thêm 7 (không phải 1). Ý 3 Sai, Ý 4 Đúng.
\end{itemize}
\begin{itemize}
    \item a) \textbf{Đúng}.
    \item b) \textbf{Đúng}.
    \item c) \textbf{Sai}.
    \item d) \textbf{Đúng}.
\end{itemize}
}
\end{ex}

\textbf{Phần 3. Trả lời ngắn}
\setcounterpageref{ex}{0}

\begin{ex} %[2D8H38]
Tìm tích của hai thừa số $2$ và $5$.
\shortans{10}
\loigiai{
\textbf{Đáp số:} $10$\\
$2 \times 5 = 10$.
}
\end{ex}

\begin{ex} %[2D8H38]
Tìm tích của hai thừa số $5$ và $8$.
\shortans{40}
\loigiai{
\textbf{Đáp số:} $40$\\
$5 \times 8 = 40$.
}
\end{ex}

\begin{ex} %[2D8V38]
Một thừa số là 2, thừa số kia là 9. Tính tích.
\shortans{18}
\loigiai{
\textbf{Đáp số:} $18$\\
$2 \times 9 = 18$.
}
\end{ex}

\begin{ex} %[2D8V38]
Một thừa số là 5, thừa số kia là 5. Tính tích.
\shortans{25}
\loigiai{
\textbf{Đáp số:} $25$\\
$5 \times 5 = 25$.
}
\end{ex}

\begin{ex} %[2D8V38]
Tính giá trị biểu thức: $2 \times 6 + 18$.
\shortans{30}
\loigiai{
\textbf{Đáp số:} $30$\\
$12 + 18 = 30$.
}
\end{ex}

\begin{ex} %[2D8C38]
Tìm tích của số 2 với số chẵn lớn nhất có một chữ số.
\shortans{16}
\loigiai{
\textbf{Đáp số:} $16$\\
Số chẵn lớn nhất có 1 chữ số là 8. Tích: $2 \times 8 = 16$.
}
\end{ex}
